\scanpage{10}

\deleted{qui est aussi la relation (unique, \inserted{à \uline{unité près}}, \unsure{que vérifie aussi bien à priori}) \unsure{vérifiée par}} Puisque $1,y,x,y^2,xy,y^3$ forment une base de $\mathcal E_6$ (adaptée à la filtration) on aura donc une équation
\[
\varphi(x,y)=\boxed{x^2+ay^3+bxy+cy^2+dx+ey+f=0}
\]
\margin{\figure{\begin{tikzpicture}\node[draw,ellipse,inner sep=3pt] {$\Delta(a,b,c,d,e,f)\in L^{\otimes12}$};\end{tikzpicture}}}
qui est d'ailleurs une équation valable dans $\mathcal E_6$ considéré comme \unsure{composante de} $\Gamma_*(f,\underline L^{\otimes3})\simeq\coprod_{i\geq0}\mathcal E_i$, satisfaite entre les sections $x,y$ de $\mathcal E_3$ \deleted{\unsure{de même $\mathcal E_3$, $y^3$}}, \inserted{à valeurs en $\mathcal E_9$, \deleted{\unsure{$\mathcal E_3$}}}.

\deleted{Cette équation} Si on désigne par $z$ la section $1$ de $\mathcal E_0=\mathcal E_1\subset \mathcal E_3$, et qu'on \deleted{\unsure{la}} écrit \inserted{notre équation} sous la forme
\[
\Phi(x,y,z)=x^2z+ay^3\deleted{\unsure{z}}+bxyz+cy^2z+dxz^2+fz^3=0
\]
on \uline{constate alors} que $\Phi$ est la forme de degré $3$ qui définit \deleted{\unsure{notre}} $C$ dans $P(\mathcal E_3)\simeq P_S^3$. En effet si $\Phi'$ est la forme qui définit $C$, alors $\Phi$ est \uline{un multiple de $\Phi'$}, et \deleted{\unsure{d'une}} c'est un multiple par quelque chose qui est inversible, puisque $\Phi$ est ``non nulle'' en tt pt de $S$ --- un de ses coefficients est $1$ ---.
\margin{\unsure{d'ailleurs, est essentiellement primitive}}
\margin{préciser ce pt cf. intersections complètes}
\mspage{9}
